\par\bigskip
\begin{center}\pdfbookmark[2]{8. Perturbations in latitude}{ceres-008}[8.]\end{center}
\noindent\hspace*{\parindent}[The formulas for the perturbations in latitude were derived similarly. Note 2 used the unperturbed orbital plane as the fixed fundamental plane. Here the ecliptic served as the fixed fundamental plane. Let] \(I,I'\) be the tangents of the inclinations and \(N,N'\) the asc[ending] nodes \(\AscendingNode\), [and] set
\[
\begin{gathered}
I\sin(v'-N)-I'\sin(v'-N')=\mathfrak I\sin(v'-\mathfrak N)\quad\mathord{\raisebox{-.7ex}{.}}\\[.7ex]
[I\cos(v'-N)-I'\cos(v'-N')=\mathfrak I\cos(v'-\mathfrak N)]
\end{gathered}
\]
then:

\SourceLeaf{401}XI) Secular motion of the node
\[
=-\tfrac14antQ'\frac{\mathfrak I}{I}\cos(N-\mathfrak N).
\]

XII) Secular increase in inclination
\[
=-\tfrac14antQ'\mathfrak I\sin(N-\mathfrak N).
\]

XIII) Periodic equations of latitude
\[
\begin{aligned}
&=\frac{a\mathfrak I}{1-\theta\theta}\left(\frac a{a'a'}-\tfrac12Q^{(0)}\right)\sin(v'-\mathfrak N)\\[.7ex]
&\quad+\frac{a\mathfrak I}{2(1-\theta)}\left\{\begin{aligned}
&\sum_{2}^{\infty}\frac{Q^{(i)}}{i+1-(i-1)\theta}\frac{\sin(iD+v'-\mathfrak N)}{i-1}\\[.7ex]
&-\sum_{1}^{\infty}\frac{Q^{(i)}}{i-1-(i+1)\theta}\frac{\sin(iD-v'+\mathfrak N)}{i+1}
\end{aligned}\right\}\mathpunct{\cdot}
\end{aligned}
\]

[The complete numerical results of the preceding derivations are published in \textsc{Zach}'s Monatliche Correspondenz, vol. VI, pp. 492—498.\refstepcounter{footnote}\hyperlink{ceres8-note}{\textsuperscript{*)}}]\begingroup
\makeatletter
\long\def\@makefntext#1{\noindent[*) #1}
\makeatother
\footnotetext{\hypertarget{ceres8-note}{}Vol. VI, pp. 227—230.]}
\endgroup{}
\par\medskip\centerline{\rule{.24\linewidth}{.4pt}}\medskip
